\documentclass[11pt,a4paper]{article}
\usepackage[a4paper,margin=2.2cm]{geometry}
\usepackage{fontspec}
\setmainfont{Times New Roman}
\setmonofont{Consolas}[Scale=0.86]
\usepackage{amsmath,amssymb,mathtools}
\usepackage{unicode-math}
\setmathfont{Latin Modern Math}
\usepackage{booktabs}
\usepackage{array}
\usepackage{enumitem}
\usepackage{titlesec}
\usepackage{xcolor}
\usepackage{mdframed}
\usepackage[hidelinks]{hyperref}

\definecolor{acc}{HTML}{7A2E1E}
\definecolor{dim}{HTML}{555555}
\definecolor{boxbg}{HTML}{F4F1EC}

\titleformat{\section}{\large\bfseries\color{acc}}{\thesection.}{0.5em}{}
\titlespacing{\section}{0pt}{1.1em}{0.4em}
\titleformat{\subsection}{\bfseries}{\thesubsection.}{0.5em}{}
\titlespacing{\subsection}{0pt}{0.8em}{0.25em}

\setlist[itemize]{leftmargin=1.2em,itemsep=0.15em,topsep=0.3em}
\setlist[enumerate]{leftmargin=1.4em,itemsep=0.15em,topsep=0.3em}
\renewcommand{\arraystretch}{1.15}
\setlength{\parskip}{0.45em}
\setlength{\parindent}{0pt}

\newcommand{\qq}[1]{\texttt{\small #1}}
\newcommand{\note}[1]{{\color{dim}\small #1}}

% Врезка для вывода или ключевого следствия.
\newmdenv[
  backgroundcolor=boxbg, linewidth=0pt, skipabove=0.7em, skipbelow=0.7em,
  innerleftmargin=0.9em, innerrightmargin=0.9em,
  innertopmargin=0.6em, innerbottommargin=0.6em
]{keybox}

\DeclareMathOperator*{\argmax}{arg\,max}
\DeclareMathOperator*{\argmin}{arg\,min}
\DeclareMathOperator{\KL}{KL}
\newcommand{\R}{\mathbb{R}}
\newcommand{\E}{\mathbb{E}}
\newcommand{\Var}{\operatorname{Var}}

\usepackage{url}
\begin{document}

{\Large\bfseries Выбор модели}

\vspace{-0.5em}
\note{vlm-recipes~· подзадача 1}

\section{Что на самом деле требуется}

В архитектуре проекта \textbf{Skill — это механизм подстановки в контекст},
а не способность модели. Методику работы с гипотезой модель получает через
\qq{select\_skill}, знать её заранее не обязана.

Значит, требования к модели узкие и измеримые:

\begin{enumerate}
\item точно исполняет инструкцию, поданную в контексте;
\item вовремя вызывает инструмент и вовремя \emph{перестаёт} вызывать;
\item держит формат ответа, который разбирается программно;
\item пригодна к дообучению.
\end{enumerate}

Предметные знания проверять не нужно — они приходят из Knowledge.
Ресурсы под инференс и обучение здесь тоже не рассматриваются: они
меняются от квартала к кварталу, а критерии — нет.

\section{Почему семейство Qwen}

\textbf{Лицензия Apache 2.0.} Без ограничений на дообучение, внедрение
и распространение производных. У Llama и Gemma лицензии с оговорками,
которые для внедрения существенны.

\textbf{Вызов инструментов обучен, а не имитируется.} В \qq{chat\_template}
есть ветка по \qq{tools}: описания функций подставляются в том формате,
в котором модель их видела при обучении. У моделей без такой ветки вызов
приходится просить текстом и разбирать вручную, надёжность заметно ниже.

\textbf{Единая линейка.} Текстовые и мультимодальные модели одного
поколения имеют общий шаблон диалога и формат данных. Переход между
размерами — замена одной строки; переход в другое семейство —
переписывание коллатора и перепроверка всего заново.

\textbf{Русский язык.} Токенизация экономнее, чем у Llama, качество выше
среднего по открытым моделям. На длинном контексте экономия токенов
переходит прямо в память и деньги.

\section{Размер: что он решает, а что нет}

Размер задаёт \emph{потолок} качества. Дообучение определяет, насколько
модель к этому потолку приблизится \emph{на вашей задаче}. Отсюда правило,
которое стоит принять до сравнения моделей:

\begin{keybox}
Дообученная модель на 7--9B на узкой задаче обгоняет модель на 70B без
дообучения — почти всегда для правил поведения и формата, почти никогда
для сложного многошагового рассуждения.
\end{keybox}

\begin{tabular}{@{}lp{10.6cm}@{}}
\toprule
\textbf{Класс} & \textbf{Роль} \\
\midrule
3--4B & судья в гардрейлах, классификация, отладка пайплайна. Не для
основной роли: держит инструкцию хуже, чем нужно агенту \\[0.3em]
7--9B & рабочая модель для узкой предметной задачи после дообучения.
Оптимум по соотношению качества к стоимости итерации \\[0.3em]
27--32B & заметный прирост на рассуждении и длинном контексте. Имеет
смысл, когда 9B упёрлась в потолок качества, а не раньше \\[0.3em]
70B и выше & когда качество решает, а ресурсы есть. Дообучение только
распределённое \\
\bottomrule
\end{tabular}

Практический порядок: отладить пайплайн и данные на 3--4B, вести
основную работу на 7--9B, подниматься выше только по результатам замера,
а не априори.

\section{Где смотреть: бенчмарки и дашборды}

Общий балл на лидерборде почти бесполезен: он усредняет несравнимое.
Для агента важны четыре узких свойства, и у каждого свой источник.

\begin{tabular}{@{}p{4.6cm}p{9.4cm}@{}}
\toprule
\textbf{Свойство} & \textbf{Где мерить} \\
\midrule
вызывает инструмент, когда нужно & BFCL, категории simple, multiple,
parallel \\[0.3em]
\textbf{не} вызывает, когда не нужно & BFCL, категория \emph{irrelevance
detection} — самая недооценённая \\[0.3em]
держит правила на длинной траектории & $\tau$-bench, метрика pass\textasciicircum k \\[0.3em]
удерживает формат ответа & IFEval, режим strict \\
\bottomrule
\end{tabular}

\subsection{Вызов инструментов}

\textbf{BFCL — Berkeley Function Calling Leaderboard.}
\url{https://gorilla.cs.berkeley.edu/leaderboard.html}\\
Основной бенчмарк по вызову функций. Смотрите разбивку по категориям,
а не итог: модели с одинаковым общим баллом различаются на irrelevance
в полтора-два раза. Многоходовые сценарии — с версии v3.

\textbf{$\tau$-bench.}
\url{https://github.com/sierra-research/tau-bench}\\
Агент в предметной области с инструментами и политикой, которую нельзя
нарушать под давлением пользователя. Ближе всего к постановке
«student-in-the-loop с Policy». Метрика pass\textasciicircum k измеряет
надёжность, а не удачу: смотрите k~=~4, а не 1.

\textbf{ToolBench}, \url{https://github.com/OpenBMB/ToolBench};
\textbf{API-Bank};
\textbf{NexusRaven}, \url{https://huggingface.co/Nexusflow}.\\
Старше и частично насыщены, но полезны как второе мнение по формату
вызовов и вложенным вызовам.

\subsection{Следование инструкциям}

\textbf{IFEval.} Инструкции с механически проверяемыми ограничениями:
«ровно три абзаца», «ответ в JSON». Без судьи — проверка регулярным
выражением; так же устроен наш \qq{skills.jsonl}. Входит в Open LLM
Leaderboard. Статья: \url{https://arxiv.org/abs/2311.07911}.

\subsection{Агентские среды}

\textbf{AgentBench}, \url{https://github.com/THUDM/AgentBench} — восемь
сред, хорош для первого отсева.
\textbf{GAIA}, \url{https://huggingface.co/spaces/gaia-benchmark/leaderboard}
— вопросы с поиском и инструментами, три уровня сложности.

\subsection{Мультимодальные модели}

\textbf{OpenVLM Leaderboard.}
\url{https://huggingface.co/spaces/opencompass/open_vlm_leaderboard}\\
Для документов смотрите DocVQA, ChartQA, OCRBench — не MMMU, который
про общие знания.

\subsection{Сводные дашборды}

\begin{tabular}{@{}lp{6.2cm}p{4.4cm}@{}}
\toprule
\textbf{Дашборд} & \textbf{Ссылка} & \textbf{Чему верить} \\
\midrule
Open LLM Leaderboard & \small\url{huggingface.co/spaces/open-llm-leaderboard/open_llm_leaderboard} & открытые модели, воспроизводимо, есть IFEval \\[0.3em]
LMArena & \small\url{lmarena.ai} & качество диалога по оценкам людей, не инструменты \\[0.3em]
Artificial Analysis & \small\url{artificialanalysis.ai} & скорость и цена по семействам \\[0.3em]
LiveBench & \small\url{livebench.ai} & обновляемые задания, меньше утечки \\[0.3em]
OpenCompass & \small\url{rank.opencompass.org.cn} & подробная разбивка по Qwen \\
\bottomrule
\end{tabular}

\vspace{0.5em}
\note{Адреса по состоянию на середину 2026 года. Лидерборды переезжают,
проверяйте.}

\section{Кто хорош в инструментах}

Семейства, устойчиво занимающие верх BFCL среди открытых моделей
по состоянию на середину 2026 года. Конкретные версии меняются каждые
несколько месяцев, поэтому здесь семейства, а не модели.

\textbf{Qwen} — верх открытой части таблицы на протяжении нескольких
поколений, вызов инструментов обучен нативно, есть все размеры.

\textbf{Llama 3.x} на 70B — сильна на BFCL, но токенизация русского
хуже, а лицензия с оговорками.

\textbf{GLM-4} и \textbf{DeepSeek-V3} — сильны в агентских сценариях,
модели преимущественно крупные.

\textbf{Mistral} — хорошая поддержка инструментов, компактные модели,
Apache 2.0.

Отдельный класс — модели, \emph{специально} дообученные под вызов
функций: \textbf{xLAM} (Salesforce), \textbf{Hermes} (Nous Research,
на базе Llama), \textbf{ToolACE}, \textbf{Functionary} (MeetKai). На BFCL
они обходят общие модели того же размера, иногда заметно.

\begin{keybox}
Но у специализированных моделей есть цена: они слабее в свободном
диалоге. Для агента, который не только вызывает инструменты, но
и разговаривает со студентом, безопаснее общая модель с нативной
поддержкой инструментов, чем узкий специалист по вызовам.
\end{keybox}

\section{Проверка кандидата}

\qq{python tools/check\_model.py <id>} — веса не качает, работает секунды.
Четыре факта, которые нужно увидеть:

\begin{tabular}{@{}ll@{}}
\toprule
\textbf{Поле} & \textbf{Требуемое значение} \\
\midrule
квантование & нет \\
вызов tools & обучен \\
русский     & не меньше 3 символов на токен \\
лицензия    & apache-2.0 \\
\bottomrule
\end{tabular}

Лидерборд отсекает заведомо слабые модели; между оставшимися выбирает
собственный замер на своих данных и своём шаблоне —
\qq{tools/selfcheck.py} и ноутбуки.

\section{Красные флаги}

\textbf{Только квантованные репозитории} (\qq{-FP8}, \qq{-NVFP4}, \qq{-AWQ},
\qq{-GPTQ}). Обучать по ним нельзя. Хуже того, \qq{transformers} не выдаёт
ошибку, а начинает разжимать веса в bf16 — обнаруживается это по
неожиданному расходу памяти. Признак в журнале: \qq{Decompressing}.

\textbf{Нет \qq{tools} в шаблоне.} Для агентского цикла модель годится
плохо: формат вызова придётся изобретать и разбирать самому.

\textbf{Неизвестный \qq{model\_type}.} Установленный \qq{transformers}
старее модели. Обновить; если поддержки нет и в основной ветке —
дообучение потребует ручной работы с кодом модели.

\section{Открытая модель против API}

Кандидаты в общей архитектуре проекта — закрытые модели по API. Для них
дообучение недоступно или ограничено, а управление поведением сводится
к системному промпту.

Смысл ветки с открытой моделью не в том, чтобы догнать их по общему
качеству, а в двух вещах, которых API не даёт: \textbf{закрепить Policy
в весах} и \textbf{вмешаться в поведение на уровне активаций}. По этой
оси её и следует оценивать.

\end{document}
